\chapter{分层习题完整解答}\label{app:answers}
本附录逐题写出计算、分类或证明步骤。读者可先独立作答，再沿运算复算；不同解法可用代回、维数、迹、行列式和等号条件互相核验。
\section{第1章：行列式}
\answer{1-1}{二阶公式给$D=2\cdot3-1\cdot5=1$。代数余子式为$C_{11}=3,C_{12}=-5,C_{21}=-1,C_{22}=2$，转置排列得
\[\adj A=\begin{pmatrix}3&-1\\-5&2\end{pmatrix}.\]
直接乘回，四个元素依次为$6-5=1,-2+2=0,15-15=0,-5+6=1$，故$A\adj A=I_2$。由于$D=1$，本题伴随矩阵也就是逆矩阵。}
\answer{1-2}{矩阵为$\begin{pmatrix}0&0&2\\0&-1&0\\3&0&0\end{pmatrix}$。按排列定义，唯一非零乘积对应列排列$(3,2,1)$，有逆序对$(3,2),(3,1),(2,1)$共3个，因此
\[D=(-1)^3\cdot2\cdot(-1)\cdot3=6.\]
也可交换第一、三行得到$\diag(3,-1,2)$，新行列式为$-6$，一次交换使原值为6。}
\answer{1-3}{转置不改行列式，$\det A\T=3$。将四阶矩阵每行都乘$-2$，四次线性性各提出一个$-2$，所以
\[\det(-2A\T)=(-2)^4\det A\T=16\cdot3=48.\]
整体倍乘的指数取矩阵阶数4，不能只提出一次$-2$。}
\answer{1-4}{设$e=(1,1,1,1)\T$，矩阵为$M=(t-1)I_4+ee\T$。在$\Span(e)$上，$ee\T e=4e$，故$M$特征值为$t+3$。在三维子空间$W=\{x:e\T x=0\}$上，$ee\T x=0$，故$M$特征值为$t-1$。任意向量都可唯一分解成这两个空间的分量，所以
\[D=(t+3)(t-1)^3.\]
全部零点是$t=-3,1$，分别重数1、3。若只用无除法更新式，$\det((t-1)I+ee\T)=(t-1)^4+4(t-1)^3$，也给同式且不遗漏$t=1$。}
\answer{1-5}{规定空行列式$D_0=1$，一阶$D_1=3$。沿末行展开：对角元3贡献$3D_{n-1}$；副对角元为1或2中的一个，其余子式再沿末列展开产生另一个，符号为负，乘积为2，故
\[D_n=3D_{n-1}-2D_{n-2}\quad(n\ge2).\]
特征方程$r^2-3r+2=(r-1)(r-2)=0$，所以$D_n=c+d2^n$。由$c+d=1,c+2d=3$解得$d=2,c=-1$，从而$D_n=2^{n+1}-1$。代入递推，右端$3(2^n-1)-2(2^{n-1}-1)=2^{n+1}-1$，核验成立。}
\answer{1-6}{按节点$0,1,2,3$的顺序，范德蒙德公式给
\[D=(1-0)(2-0)(3-0)(2-1)(3-1)(3-2)=1\cdot2\cdot3\cdot1\cdot2\cdot1=12.\]
若把节点作为行而非列，整个矩阵只是转置，值仍为12；若仅交换两节点，则值变号。}
\answer{1-7}{写$A=(a_1,\ldots,a_n)$，更新后第$j$列为$a_j+v_j u$。按列多线性展开。选取两个或更多更新列的项都含两条平行于$u$的列，行列式为零；未选更新列的项为$\det A$。恰选第$j$个更新列时，贡献为$v_j\det(a_1,\ldots,u,\ldots,a_n)$。沿第$j$列展开，此行列式为$\sum_i C_{ij}u_i$，其中余子式仍取自$A$，因为删去了被替换的那列。因此
\[\det(A+uv\T)=\det A+\sum_{j,i}v_j C_{ij}u_i=\det A+v\T\adj(A)u.\]
证明全程没有除以$\det A$，故包含奇异$A$。}
\answer{1-8}{在上一题中令$A=I$、更新向量取$tu,v$。因为$\adj I=I$，立即得
\[f(t)=\det(I+tuv\T)=1+t\,v\T u.\]
这是关于$t$的精确一次式，故$f'(t)=v\T u$，特别$f'(0)=v\T u$。当$v\T u=0$时，行列式对全部$t$恒为1，不能据此误判更新矩阵必等于$I$。}
\answer{1-9}{幂等性给$V=\im A\oplus\ker A$，在适配基下$A$为$\diag(I_r,0)$。因而$I+A$为$\diag(2I_r,I_{n-r})$，所以$\det(I+A)=2^r$，而$\tr A=r$。比较$2^r$与$1+r$：$r=0,1$时相等；$r=2$时$4>3$。若$r\ge2$且$2^r>r+1$，则$2^{r+1}>2r+2\ge r+2$，归纳可知所有$r\ge2$均严格大于。因此两式相等当且仅当$r\le1$。}
\section{第2章：矩阵}
\answer{2-1}{令$N=\begin{pmatrix}0&2\\0&0\end{pmatrix}$，则$A=I+N,N^2=0$。计算$(I+N)(I-N)=I$得$A\inv=\begin{pmatrix}1&-2\\0&1\end{pmatrix}$。二项式中二次及以上$N$项为零，故
\[A^3=I+3N+3N^2+N^3=\begin{pmatrix}1&6\\0&1\end{pmatrix}.\]
直接算$A^2=\begin{pmatrix}1&4\\0&1\end{pmatrix}$再乘$A$，也得到上式。}
\answer{2-2}{这里$A,B$都是$n$阶方阵。由$AB=I_n$得$n=\rank AB\le\rank A\le n$，故$\rank A=n$，$A$可逆。在$AB=I$两端左乘$A\inv$得$B=A\inv$，于是$BA=A\inv A=I$。也可取行列式得到$\det A\det B=1$，进而证明可逆。矩形矩阵不能照用此结论。}
\answer{2-3}{行变换$R_2\leftarrow R_2-2R_1$得到$\begin{pmatrix}1&2&3\\0&0&0\end{pmatrix}$，非零行恰一条，所以秩1。独立核验：元素1构成非零一阶子式，秩至少1；两行成比例，秩至多1，合并得1。}
\answer{2-4}{移项得$3A-A^2=2I$，即
\[A\dfrac{3I-A}{2}=I,\qquad\dfrac{3I-A}{2}A=I.\]
两侧相等因为$3I-A$是$A$的多项式，故所给表达同时为左右逆，不需预先假定$A$可逆。因此$A\inv=\frac{3I-A}{2}$。}
\answer{2-5}{秩2说明某个二阶子式非零，因此某个代数余子式非零，$\adj A\ne0$。又三阶$A$奇异，$A\adj A=(\det A)I=0$，故$\adj A$的各列都在一维核$\ker A$中，给$\rank\adj A\le1$。结合非零得$\rank\adj A=1$。三阶矩阵秩1小于3，故$\det\adj A=0$。}
\answer{2-6}{设$X=\begin{pmatrix}x_{11}&x_{12}\\x_{21}&x_{22}\end{pmatrix}$。左乘对角矩阵使第一行乘2、第二行乘3，因此$2x_{11}=2,2x_{12}=4,3x_{21}=3,3x_{22}=6$，得
\[X=\begin{pmatrix}1&2\\1&2\end{pmatrix}.\]
左系数矩阵可逆，故此解唯一。也可在方程两侧左乘$\diag(\frac{1}{2},\frac{1}{3})$一次得到它。}
\answer{2-7}{先提取标量：$2I-N=2(I-\frac{N}{2})$。记$E=\frac{N}{2}$，则$E^3=0$。有限几何和满足$(I-E)(I+E+E^2)=I-E^3=I$，左右顺序交换也成立，故
\[(2I-N)\inv=\dfrac12(I+\dfrac{N}{2}+\dfrac{N^2}{4})=\dfrac12I+\dfrac14N+\dfrac18N^2.\]
逐项乘回为$I+\frac{N}{2}+\frac{N^2}{4}-\frac{N}{2}-\frac{N^2}{4}-\frac{N^3}{8}=I$，验证逆存在。}
\answer{2-8}{设$w=A\inv u,s=v\T w$，则$A+uv\T=A(I+wv\T)$。试设$(I+wv\T)\inv=I-cwv\T$。利用$(wv\T)^2=s\,wv\T$，乘积变为$I+[1-c(1+s)]wv\T$，当$1+s\ne0$取$c=\frac{1}{1+s}$即可；相反顺序也给$I$。右乘$A\inv$得到
\[(A+uv\T)\inv=A\inv-\dfrac{A\inv uv\T A\inv}{1+v\T A\inv u}.\]
若$1+s=0$，则$w\ne0$，而$(I+wv\T)w=(1+s)w=0$，所以更新矩阵奇异。因此分母非零不仅充分，而且必要。}
\answer{2-9}{设$A$为$m\times n$、秩$r>0$。选取$r$个极大无关列$a_{j_1},\ldots,a_{j_r}$，组成$m\times r$矩阵$B$。每列$a_j$都能唯一写为$Bc_j$；按列拼接$r\times n$矩阵$C=(c_1,\ldots,c_n)$，即有$A=BC$。$B$的列本来无关，故满列秩。对于被选的第$j_k$列，其坐标$c_{j_k}=e_k$，所以$C$含$I_r$作为列子矩阵，秩至少$r$，又行数为$r$，故满行秩。$r=0$时$A=0$，可取$m\times0$和$0\times n$的空因子，其乘积约定为零矩阵。}
\section{第3章：向量组}
\answer{3-1}{记$a=(1,2)\T,b=(2,4)\T$，则$2a-b=0$，系数$(2,-1)$不全零，所以二者相关。$a\ne0$，单向量组$(a)$无关，又$b=2a$，所以它是极大无关组，组秩1；也可取$(b)$。}
\answer{3-2}{记题给列为$a_1=e_1,a_2=e_2,a_3=e_1+e_2$。若$c_1a_1+c_2a_2=0$，比较前两分量得$c_1=c_2=0$，故前两列无关。第三列为其和，故$(a_1,a_2)$生成原组，成为极大无关组，秩2。其他两组$(a_1,a_3),(a_2,a_3)$也无关且生成原组，体现极大无关组不唯一。}
\answer{3-3}{设坐标为$(c_1,c_2)\T$，由
\[(4,2)\T=c_1(1,1)\T+c_2(1,-1)\T\]
得$c_1+c_2=4,c_1-c_2=2$。两式相加得$2c_1=6$，故$c_1=3,c_2=1$。所求坐标为$(3,1)\T$；代回$3(1,1)+(1,-1)=(4,2)$核验。}
\answer{3-4}{设$u=a+tb,v=2a+b$。若$\alpha u+\beta v=0$，按原无关组$a,b$整理得$(\alpha+2\beta)a+(t\alpha+\beta)b=0$，所以
\[\begin{pmatrix}1&2\\t&1\end{pmatrix}\begin{pmatrix}\alpha\\\beta\end{pmatrix}=0.\]
系数行列式为$1-2t$，故$t\ne\frac{1}{2}$时仅有零解，新组无关。$t=\frac{1}{2}$时$v=2u$，且$u\ne0$，所以相关、秩1。这一分支必须单列，不能在除以$1-2t$后丢掉。}
\answer{3-5}{令$P=\begin{pmatrix}2&1\\0&1\end{pmatrix}$，因$C=BP$，同一向量的坐标满足$[x]_B=P[x]_C$。设新坐标$(u,v)\T$，得到$2u+v=5,v=1$，故$u=2$。也可算$P\inv=\begin{pmatrix}\frac{1}{2}&-\frac{1}{2}\\0&1\end{pmatrix}$，乘$(5,1)\T$得$(2,1)\T$。把它再乘$P$返回$(5,1)\T$，检查方向无误。}
\answer{3-6}{两向量$a=(1,1,1)\T,b=(1,0,-1)\T$内积为$1+0-1=0$，所以Schmidt第二步减去的投影为零。长度$\norm a=\sqrt3,\norm b=\sqrt2$，故单位正交组为
\[q_1=\dfrac1{\sqrt3}(1,1,1)\T,\qquad q_2=\dfrac1{\sqrt2}(1,0,-1)\T.\]
各自长度为1、相互内积0，并且只对原列作非零倍乘，所以生成空间不变。}
\answer{3-7}{设$v_1,\ldots,v_k$无关，$T$可逆。若$\sum_i c_iT(v_i)=0$，由线性性有$T(\sum_i c_iv_i)=0$；施加$T\inv$得$\sum_i c_iv_i=0$。原组无关迫使全部$c_i=0$，故像组无关。实际上只需$\ker T=\{0\}$即单射；若无此条件，非零核向量可能被送成零，结论失败。}
\answer{3-8}{记$q_1=1,q_2=1+x,q_3=1+x+x^2$，有显式表示$1=q_1,x=q_2-q_1,x^2=q_3-q_2$，故它们生成$\mathcal P_2$。再验证无关：若$u q_1+v q_2+w q_3=0$，比较$x^2,x,1$系数依次得$w=0,v+w=0,u+v+w=0$，所以$u=v=w=0$。故$q_i$构成基，每个多项式表示唯一；一般$a+bx+cx^2=(a-b)q_1+(b-c)q_2+cq_3$。}
\answer{3-9}{在实内积空间中展开
\[\begin{aligned}
\norm{a+b}^2&=\norm a^2+2a\T b+\norm b^2,\\
\norm{a-b}^2&=\norm a^2-2a\T b+\norm b^2.
\end{aligned}\]
两式相等即$4a\T b=0$，故$a\perp b$。反过来内积为零也使两长度相等，几何上对应平行四边形两条对角线等长。该写法默认实内积；复内积时相等只推出实部为零。}
\section{第4章：线性方程组}
\answer{4-1}{\textbf{消元法：}第二式给$y=x-1$，代入第一式得$2x+x-1=5$，所以$x=2,y=1$。代回$2\cdot2+1=5,2-1=1$。
\textbf{克拉默法：}系数矩阵$A=\begin{pmatrix}2&1\\1&-1\end{pmatrix}$，$D=-2-1=-3\ne0$。替换第一列得$D_x=\det\begin{pmatrix}5&1\\1&-1\end{pmatrix}=-6$，替换第二列得$D_y=\det\begin{pmatrix}2&5\\1&1\end{pmatrix}=-3$，故$x=\frac{D_x}{D}=2,y=\frac{D_y}{D}=1$。非零$D$也保证唯一性。}
\answer{4-2}{系数秩为1，未知量3个，因此核维数$3-1=2$。设$y=s,z=t$，则$x=-2s-t$，所以
\[\vx=s(-2,1,0)\T+t(-1,0,1)\T.\]
两列代入原齐次式都为零。若它们以系数$u,v$组合为零，比较第二、三分量得$u=v=0$，故无关；又已表示全部解，所以构成基础解系。}
\answer{4-3}{先取特解$\vx_p=(3,0,0)\T$，它满足$x+2y+z=3$。任意另一解$\vx$与它相减满足齐次式，故由前题
\[\vx=(3,0,0)\T+s(-2,1,0)\T+t(-1,0,1)\T,\quad s,t\in\R.\]
反过来代入该表达，左端恒为$3+s(-2+2)+t(-1+1)=3$，所以表达既不漏解，也不添解。}
\answer{4-4}{用第二式减第一式，得到$(t-1)x=1$。当$t=1$时左端恒为0，与右端1矛盾，无解。当$t\ne1$时
\[x=\dfrac1{t-1},\qquad y=1-x=\dfrac{t-2}{t-1}.\]
此时系数行列式$1-t\ne0$，故解唯一。检查第二式左端为$\frac{t+(t-2)}{t-1}=2$。所有参数已分完，没有无穷多解的分支。}
\answer{4-5}{秩--零化度给齐次解空间维数$4-2=2$。相容非齐次解集为$\vx_p+\ker A$，因此也有两个自由参数。若右端$\vb\ne0$，零向量不满足$A0=\vb$，所以解集不是线性空间，而是二维仿射平移；若右端实际为零，则它本身就是二维线性空间。“相容”保证特解存在，不能省掉这个前提。}
\answer{4-6}{由线性性以及$A\eta_1=A\eta_2=\vb$，有
\[A(2\eta_1-\eta_2)=2\vb-\vb=\vb,\qquad A(\eta_1+\eta_2)=2\vb.\]
因此第一向量仍解原系统，第二向量解右端$2\vb$的系统。题设$\vb\ne0$确保$2\vb\ne\vb$，所以第二个不解原系统。一般$\sum c_i\eta_i$的右端为$(\sum c_i)\vb$，系数和为1时保留原右端。}
\answer{4-7}{设$A$为实$m\times n$且满列秩。任意$\vx\ne0$有$A\vx\ne0$，所以$\vx\T A\T A\vx=\norm{A\vx}^2>0$。由此$A\T A$正定，特别可逆。定义
\[L=(A\T A)\inv A\T\in\R^{n\times m},\]
直接相乘$LA=(A\T A)\inv(A\T A)=I_n$，即得到左逆。一般$AL$是到$\im A$的投影，不要求等于$I_m$。}
\answer{4-8}{三个方程为$x_1=a,x_2=b,x_1+x_2=c$，前两式已唯一规定未知量，所以相容当且仅当$c=a+b$；成立时解为$(a,b)\T$。另可取左零向量$y=(-1,-1,1)\T$，它满足$A\T y=0$，相容必有$y\T(a,b,c)\T=-a-b+c=0$。由于$A$秩2，左核维数$3-2=1$，此向量张成全部左核，因此该条件也充分。}
\answer{4-9}{\textbf{不等式法：}Cauchy不等式给$9=(x+y+z)^2\le(1^2+1^2+1^2)(x^2+y^2+z^2)$，故范数平方至少3。等号当且仅当$(x,y,z)$与$(1,1,1)$成比例，结合分量和3得$(1,1,1)$。
\textbf{正交分解法：}任意解写成$(1,1,1)+h$，其中$h_1+h_2+h_3=0$，所以$h\perp(1,1,1)$，从而$\norm\vx^2=3+\norm h^2$。唯一最优是$h=0$，最小范数为$\sqrt3$。这同时说明最小范数解位于齐次核的正交补。}
\section{第5章：特征值与相似}
\answer{5-1}{三角矩阵的特征多项式为$(\lambda-1)(\lambda-3)$，故根为1、3。对$\lambda=1$，
\[A-I=\begin{pmatrix}0&2\\0&2\end{pmatrix}\]
给$y=0,x$自由，所以特征空间为$\Span((1,0)\T)$。对$\lambda=3$，$A-3I=\begin{pmatrix}-2&2\\0&0\end{pmatrix}$给$x=y$，特征空间为$\Span((1,1)\T)$。各空间的非零向量才是特征向量，零向量本身不算。}
\answer{5-2}{按代数重数求和、求积，得$\tr A=-1+2+2=3$，$\det A=(-1)\cdot2\cdot2=-4$。即使$A$不可对角化，也能在复数域相似上三角化，其平方的对角为各特征值平方，故$\tr(A^2)=(-1)^2+2^2+2^2=9$。不能用$(\tr A)^2$冒充$\tr(A^2)$；本题两者恰巧同为9，只是数值巧合。}
\answer{5-3}{给定矩阵已经对角，取$P=I_3$即可对角化。根2的特征空间由$(A-2I)x=0$得到$x_3=0$，为$\Span(e_1,e_2)$，几何重数2等于代数重数2；根3的空间为$\Span(e_3)$。三条独立特征向量共同构成基，因此重根不妨碍对角化。}
\answer{5-4}{特征多项式恒为$(\lambda-1)^2$。若$t=0$，$A=I_2$，特征空间维数2，可以对角化。若$t\ne0$，$A-I=\begin{pmatrix}0&t\\0&0\end{pmatrix}$秩1，故根1的特征空间维数$2-1=1$，小于代数重数2，不能对角化。等价地，只有一个特征值的可对角化矩阵必为相应标量矩阵，而$t\ne0$时它不是。}
\answer{5-5}{特征方程$(3-\lambda)^2-1=0$给$\lambda=4,2$。根4满足$y=x$，根2满足$y=-x$，取单位特征向量$q_1=\frac{(1,1)\T}{\sqrt2},q_2=\frac{(1,-1)\T}{\sqrt2}$，两者点积为0。令
\[Q=\dfrac1{\sqrt2}\begin{pmatrix}1&1\\1&-1\end{pmatrix}.\]
则$Q\T Q=I$，由$Aq_1=4q_1,Aq_2=2q_2$得$AQ=Q\diag(4,2)$，左乘$Q\T$即$Q\T AQ=\diag(4,2)$。列的次序与对角特征值次序必须配套。}
\answer{5-6}{设$B=P\inv AP$，则$B-\lambda I=P\inv(A-\lambda I)P$。当$k=1$公式成立；若$k$时成立，乘一次并约去中间$PP\inv$，得
\[(B-\lambda I)^{k+1}=P\inv(A-\lambda I)^kPP\inv(A-\lambda I)P=P\inv(A-\lambda I)^{k+1}P.\]
故所有非负整数$k$成立，$k=0$时两边为$I$。可逆左、右乘保秩，所以两次幂秩相同。比较的是“先取幂再求秩”，而非秩这个整数的$k$次幂。}
\answer{5-7}{由$A^2=A$，任意向量有$x=Ax+(I-A)x$，分别落在$\im A$与$\ker A$。交中向量$w=Au$又满足$Aw=0$，但$Aw=A^2u=Au=w$，所以交为零。选两空间的基拼接，$A$在像空间上恒等、核上为零，故相似于$\diag(I_r,0)$，从而可对角化且谱只能含0、1。另由特征对$Av=\lambda v$直接得$\lambda^2v=\lambda v$，也说明$\lambda\in\{0,1\}$。}
\answer{5-8}{令$J_3(1)=I+N$，其中$N$第一上副对角全为1，$N^2$只有$(1,3)$位置为1，$N^3=0$。对非负整数$m$，可交换二项式展开给
\[(I+N)^m=I+mN+\binom m2N^2=\begin{pmatrix}1&m&\dfrac{m(m-1)}{2}\\0&1&m\\0&0&1\end{pmatrix}.\]
因此对角元全1，第一上副对角为$m$，第二上副对角为$\frac{m(m-1)}{2}$；$m=0,1$时某些上三角项为零。若题意包含负整数，公式仍成立：先有$(I+N)^{-1}=I-N+N^2$，并且上式右端的矩阵$G(m)$满足$G(m)G(\ell)=G(m+\ell)$，由$\binom{m+\ell}2=\binom m2+\binom\ell2+m\ell$验证即可。}
\answer{5-9}{实对称谱定理给$A=Q\Lambda Q\T$，$Q$正交。置$y=Q\T x$，则$\norm y=\norm x$，且
\[x\T Ax=\sum_i\lambda_i y_i^2.\]
每项$y_i^2\ge0$，由$2\le\lambda_i\le5$分别乘后相加，得到$2\sum_i y_i^2\le\sum_i\lambda_i y_i^2\le5\sum_i y_i^2$，即所求。下界等号要求全部非零坐标都对应特征值2；上界同理对应5。若相应端点不在谱中，则非零$x$不会取到该端点等号。}
\section{第6章：二次型}
\answer{6-1}{实对称矩阵$A=\begin{pmatrix}a&b\\b&c\end{pmatrix}$对应$x\T Ax=ax^2+2bxy+cy^2$。比较系数$a=1,2b=6,c=4$，故
\[A=\begin{pmatrix}1&3\\3&4\end{pmatrix}.\]
交叉项由两个对称位置共同贡献，因此每处为系数的一半。}
\answer{6-2}{令$u=\sqrt2x,v=\sqrt3y$，此线性代换行列式$\sqrt6\ne0$，可逆。二次型成为$u^2-v^2$，即规范形。正指数$p=1$，负指数$q=1$，零指数0，秩$p+q=2$；能取正负两种值，所以是不定型。原系数2、$-3$不必保留在规范形中，惯性只记符号个数。}
\answer{6-3}{二次型为$x_1^2+2x_3^2$，任意实向量都使它非负，所以半正定。非零向量$e_2$使值为零，故非正定。矩阵仅第一、三列非零且独立，秩2；惯性为$(2,0)$，零指数1。特征值$1,0,2$也给同样判断。}
\answer{6-4}{\textbf{配方法：}
\[3x^2+2xy+3y^2=3(x+\dfrac{y}{3})^2+(3-\dfrac{1}{3})y^2=3(x+\dfrac{y}{3})^2+\dfrac83y^2.\]
取$u=x+\frac{y}{3},v=y$，逆变换$x=u-\frac{v}{3},y=v$，行列式为1，得到$3u^2+\frac{8v^2}{3}$。
\textbf{正交法：}矩阵$\begin{pmatrix}3&1\\1&3\end{pmatrix}$的特征值为4、2，单位特征向量为$\frac{(1,1)\T}{\sqrt2},\frac{(1,-1)\T}{\sqrt2}$。取$x=\frac{s+t}{\sqrt2},y=\frac{s-t}{\sqrt2}$，则$x^2+y^2=s^2+t^2,2xy=s^2-t^2$，二次型成为$4s^2+2t^2$。两结果的对角数不同，但正指数均2，均可再缩放成规范形$u_1^2+u_2^2$。}
\answer{6-5}{顺序主子式为$\Delta_1=t,\Delta_2=3t-4$。正定充要条件是二者严格正，合并为$t>\frac{4}{3}$。也可完成平方：$tx^2+4xy+3y^2=3(y+\frac{2x}{3})^2+(t-\frac{4}{3})x^2$，这直接给同一条件。边界$t=\frac{4}{3}$只剩一个平方，是秩1半正定；$t<\frac{4}{3}$时取$y=-\frac{2x}{3},x\ne0$得到负值，故不正定。}
\answer{6-6}{对角型$-x_1^2-2x_2^2+3x_3^2$有一个正项、两个负项，所以惯性$(p,q)=(1,2)$，秩3。直线$\Span(e_3)$上任意非零向量使值为正，故正定子空间至少一维。若存在二维正定子空间$L$，它与二维负子空间$N=\Span(e_1,e_2)$在三维空间中的交维数至少$2+2-3=1$，交中非零向量会同时使二次型正、负，矛盾。故最大维数恰1。}
\answer{6-7}{谱定理给$A=\sum_{i=1}^n\lambda_iq_iq_i\T$，其中$q_i$单位正交。半正定使$\lambda_i\ge0$；秩$r$说明恰有$r$个正特征值，将它们排在前面。取$u_i=\sqrt{\lambda_i}q_i$，则每个$u_i\ne0$且
\[A=\sum_{i=1}^r u_iu_i\T.\]
这也证明$r$项足够。反之每个非零外积秩1，若能用$k$项表示，则秩和不等式给$r=\rank A\le k$，所以不能少于$r$项。$A=0$时$r=0$，用空和表示。}
\answer{6-8}{由约束，二次型为$3+2xy$。不等式$(x-y)^2\ge0$给$2xy\le x^2+y^2=1$，$(x+y)^2\ge0$给$2xy\ge-1$，所以值在$[2,4]$。上界在$x=y=\pm\frac{1}{\sqrt2}$取到，下界在$x=-y=\pm\frac{1}{\sqrt2}$取到。谱法中矩阵特征值4、2分别对应和方向与差方向，单位特征向量恰为上述点，所以几何上是单位圆上的主轴极值。}
\answer{6-9}{正定$A$可逆，行列式秩一更新式给
\[\det(A+uu\T)=\det A\,(1+u\T A\inv u).\]
由谱分解$A=Q\diag(\lambda_i)Q\T$、$\lambda_i>0$，有$\det A=\prod_i\lambda_i>0$，且$A\inv=Q\diag(\frac{1}{\lambda_i})Q\T\succ0$。非零$u$使$u\T A\inv u>0$，所以括号严格大于1，再乘正数$\det A$即得严格不等式。若$u=0$则等号成立，说明非零条件不可删。}
\section{第7章：线性空间}
\answer{7-1}{集合含零多项式。若$p(0)=q(0)=0$，则任意标量$\alpha,\beta$有$(\alpha p+\beta q)(0)=0$，所以是子空间。设$p=a+bx+cx^2$，条件$p(0)=0$给$a=0$，故$p=bx+cx^2$。$x,x^2$生成该空间，且$ux+vx^2=0$比较系数得$u=v=0$，因此它们是基，维数2。}
\answer{7-2}{设$K\T=-K$。对角元素满足$k_{ii}=-k_{ii}$，实数域中得$k_{ii}=0$；下三角由上三角负值决定。因此
\[K=a(E_{12}-E_{21})+b(E_{13}-E_{31})+c(E_{23}-E_{32}).\]
三矩阵生成全部反对称矩阵；线性组合为零时分别比较$(1,2),(1,3),(2,3)$元素得$a=b=c=0$，故为基，维数3。}
\answer{7-3}{核满足$x+y=0,y+z=0$。令$y=s$，得$(x,y,z)=s(-1,1,-1)$，故核维数1。映射矩阵$A=\begin{pmatrix}1&1&0\\0&1&1\end{pmatrix}$含由第一、三列组成的$I_2$，故秩2。更直接地，对任意$(a,b)\in\R^2$，有$T(a,0,b)=(a,b)$，所以像为$\R^2$，维数2；两维数相加为3。}
\answer{7-4}{$U$中向量形如$(a,b,0)$，$W$中形如$(0,c,d)$，同时属于两者要求第一、三分量为零，所以$U\cap W=\Span(e_2)$，维数1。和包含$e_1,e_2,e_3$，故$U+W=\R^3$，维数3。也可写任意$(x,y,z)=(x,y,0)+(0,0,z)$证生成。因为交非零，这个和不是直和；公式$2+2-1=3$核验维数。}
\answer{7-5}{令$a=(1,1)\T,b=(2,1)\T$。投影系数为$c=\frac{a\T b}{a\T a}=\frac{3}{2}$，所以$p=ca=(\frac{3}{2},\frac{3}{2})\T$，残差$r=b-p=(\frac{1}{2},-\frac{1}{2})\T$。验证$a\T r=0$，故残差垂直投影直线。对直线上任意点$sa$，勾股分解$\norm{b-sa}^2=\norm r^2+\norm{(c-s)a}^2$说明$c=\frac{3}{2}$给唯一最近点。}
\answer{7-6}{目标平方和为
\[f(c)=(c-1)^2+(c-2)^2+(c-6)^2=3c^2-18c+41=3(c-3)^2+14.\]
故唯一最小值14在$c=3$取得。矩阵法取$A=(1,1,1)\T,\vb=(1,2,6)\T$，正规方程为$(A\T A)c=A\T\vb$，即$3c=9$，也得$c=3$。残差$(-2,-1,3)\T$分量和零，因而与常数列正交。}
\answer{7-7}{设$\dim V=n<\infty$。单射等价于$\ker T=\{0\}$，秩--零化度给$\dim\im T=n-0=n$。像是$V$的$n$维子空间，故只能等于$V$，即满射。无限维反例取实多项式空间$\R[x]$，令$T(p)=xp$。若$T(p)=0$，比较系数知$p=0$，所以单射；但任何$xp$的常数项为零，不能等于1，因此不满射。有限维保证维数计数能推出相等，反例说明不可直接删掉此条件。}
\answer{7-8}{任何$p=a+bx+cx^2+dx^3$的陪集满足$p+\mathcal P_1=c(x^2+\mathcal P_1)+d(x^3+\mathcal P_1)$，所以这两个陪集生成商空间。若其以系数$u,v$组合为零陪集，则$ux^2+vx^3\in\mathcal P_1$，比较二、三次项得$u=v=0$。因此它们构成基，商维数2，也等于$\dim\mathcal P_3-\dim\mathcal P_1=4-2$。}
\answer{7-9}{\textbf{必要性：}若$P$是到子空间$W$的正交投影，取$W$单位正交基列成$Q$，则$Px=\sum_i q_i(q_i\T x)$，故$P=QQ\T$。于是$P\T=P$，$P^2=Q(Q\T Q)Q\T=QQ\T=P$。
\textbf{充分性：}设$P\T=P,P^2=P$。对任意$x$，$x=Px+(I-P)x$，且$P(I-P)x=0$，所以两项分别在$\im P,\ker P$。若$u=Py,v\in\ker P$，则$u\T v=y\T P\T v=y\T Pv=0$，两空间正交。因此$Px$正是$x$在$\im P$上的正交分量，$P$是正交投影。}
\section{第8章：提高专题}
\answer{8-1}{Sylvester不等式中公共维数为5，故$\rank AB\ge4+3-5=2$；乘积秩不超过任一因子，故$\rank AB\le\min(4,3)=3$。两界均可达到：取$A=\diag(1,1,1,1,0)$。取$B=\diag(1,1,0,0,1)$时$AB$秩2；取$B=\diag(1,1,1,0,0)$时$AB$秩3。两种$B$的秩都是3，所以范围不能进一步普遍缩小。}
\answer{8-2}{$I$与$2N$可交换，用二项式定理展开。由于$N^2=0$，所有二次及以上项为零，故
\[(I+2N)^5=I+\binom51 2N+\binom52(2N)^2+\cdots+(2N)^5=I+10N.\]
也可反复乘法归纳$(I+2N)^k=I+2kN$，在$k=5$得到同式。}
\answer{8-3}{对角矩阵上$p(A)=\diag(p(1),p(1),p(2))$。因此零化的充要条件是$p(1)=p(2)=0$，即互异因子$t-1,t-2$都整除$p$，从而其乘积整除$p$。该乘积本身零化$A$，又是首一，所以最小多项式为$(t-1)(t-2)$。根1重复出现两次只增加特征多项式重数，不增加最小多项式指数。}
\answer{8-4}{计算$\chi_A(t)=\det\begin{pmatrix}t-2&-1\\-1&t-1\end{pmatrix}=(t-2)(t-1)-1=t^2-3t+1$。CH给$A^2-3A+I=0$，即$A(3I-A)=(3I-A)A=I$，故
\[A\inv=3I-A=\begin{pmatrix}1&-1\\-1&2\end{pmatrix}.\]
相乘得到$\begin{pmatrix}2-1&-2+2\\1-1&-1+2\end{pmatrix}=I$，验证符号与常数无误。}
\answer{8-5}{四阶矩阵核维数为$d_k=4-\rank N^k$，所以$d_0=0,d_1=2,d_2=3,d_3=4,d_4=4$。连续差为$2,1,1,0$，分别数出大小至少1、2、3、4的块数。因此恰为大小1的有$2-1=1$个，恰为大小2的有$1-1=0$个，恰为大小3的有$1-0=1$个。故Jordan形为$J_3(0)\oplus J_1(0)$，其秩依次为2、1、0，与题给相符；最小多项式为$t^3$。}
\answer{8-6}{矩阵1范数取最大列绝对值和，两列分别为$|1|+|3|=4$、$|-2|+|4|=6$，故$\norm A_1=6$。无穷范数取最大行绝对值和，两行分别为$1+2=3$、$3+4=7$，故$\norm A_\infty=7$。Frobenius范数为全部元素平方和开根，故$\norm A_F=\sqrt{1+4+9+16}=\sqrt{30}$。三种范数不应混用行和、列和或绝对值。}
\answer{8-7}{设$m_A$是首一、次数最小的非零零化多项式。任意非零$p$满足$p(A)=0$，用多项式带余除法写$p=qm_A+r$，其中$r=0$或$\deg r<\deg m_A$。代入矩阵得$r(A)=p(A)-q(A)m_A(A)=0$。若$r\ne0$，除以最高次系数可得到更低次首一零化多项式，违反最小性，所以$r=0$，即$m_A\mid p$；$p=0$时整除也成立。若$m_A,\widetilde m_A$都是最小多项式，则彼此整除，只能相差非零常数倍；又都首一，故完全相等。}
\answer{8-8}{可逆矩阵的全部奇异值严格正，故$\kappa_2(A)=\frac{\sigma_{\max}}{\sigma_{\min}}\ge1$。等号当且仅当最大、最小奇异值相等，即全部奇异值同为$s>0$。此时SVD给$A=U(sI)V\T=sQ$，其中$Q=UV\T$且$Q\T Q=I$，所以是正数倍正交矩阵。反过来若$A=sQ$且$Q$正交、$s\ne0$，则$A\T A=s^2I$，全部奇异值为$|s|$，条件数为1。等价条件还可写为$A\T A=s^2I$。}
\answer{8-9}{置$q=\norm E<1,S_k=\sum_{j=0}^kE^j$。诱导范数满足$\norm I=1$及次乘性，故对$\ell>k$有
\[\norm{S_\ell-S_k}\le\sum_{j=k+1}^{\ell}\norm E^j\le\sum_{j=k+1}^{\ell}q^j\le\dfrac{q^{k+1}}{1-q}\longrightarrow0.\]
有限维矩阵空间完备，所以$S_k$收敛到矩阵$S$。有限几何和恒等式给$(I-E)S_k=S_k(I-E)=I-E^{k+1}$，且$\norm{E^{k+1}}\le q^{k+1}\to0$。取极限，乘法连续性给$(I-E)S=S(I-E)=I$，因此$I-E$可逆且逆为$S$。最后
\[\norm{(I-E)\inv}=\norm S\le\sum_{j=0}^{\infty}\norm E^j\le\sum_{j=0}^{\infty}q^j=\dfrac1{1-q}.\]
若$q=0$则$E=0$，上述结论同样成立。}
